% chktex-file 44
% chktex-file 1

\subsection{Pengujian Fungsional}

Pengujian fungsional memverifikasi alur pengguna dan pipeline F01--F10. Satu skenario dapat mencakup beberapa kebutuhan yang terjadi dalam satu alur, tetapi seluruh kebutuhan memiliki paling sedikit satu bukti pengujian pada Tabel~\ref{tab:skenario-fungsional}.

\begin{styledxltabular}{|>{\raggedright\arraybackslash}p{0.09\textwidth}|>{\raggedright\arraybackslash}p{0.13\textwidth}|>{\raggedright\arraybackslash}X|>{\raggedright\arraybackslash}X|}
    \hiderowcolors
    \caption{Skenario Pengujian Fungsional Sistem DecMed-PGHD}\label{tab:skenario-fungsional} \\
    \hline
    \rowcolor{headerBg}
    \textbf{Kode} & \textbf{Kebutuhan} & \textbf{Skenario Pengujian} & \textbf{Hasil yang Diharapkan} \\
    \hline
    \showrowcolors
    \endfirsthead
    \hline
    \rowcolor{headerBg}
    \textbf{Kode} & \textbf{Kebutuhan} & \textbf{Skenario Pengujian} & \textbf{Hasil yang Diharapkan} \\
    \hline
    \endhead
    \hline
    \endfoot
    \hline
    \endlastfoot
    T-F01 & F01 & Pasien memberi izin Health Connect/sensor, menjalankan pengumpulan otomatis, kemudian menambahkan satu record manual. & Record otomatis dan manual tersimpan dengan nilai, waktu pengukuran, dan sumber yang sesuai. \\
    \hline
    T-F02 & F02 & Perangkat dibuat luring, PGHD dikumpulkan, lalu aplikasi ditutup dan dibuka kembali. & Record tetap tersedia pada basis data lokal dan belum ditandai sebagai terkirim. \\
    \hline
    T-F03 & F03 & Beberapa record yang belum dikirim dibentuk menjadi batch ketika perangkat luring, kemudian koneksi diaktifkan kembali. & Batch tetap tersimpan saat luring dan pengiriman dapat dijalankan ketika koneksi tersedia. \\
    \hline
    T-F04 & F04 & Pasien membuka halaman record dan batch setelah proses koleksi serta pengiriman. & Aplikasi menampilkan sumber, waktu, jumlah record, identitas batch, dan status pengiriman. \\
    \hline
    T-F05 & F05 & PGHD dikumpulkan dari minimal dua sumber, misalnya Health Connect dan input manual, lalu isi payload diperiksa. & Payload mempertahankan identitas pasien, waktu, sumber, perangkat/aplikasi, metode pencatatan, dan status validasi yang tersedia. \\
    \hline
    T-F06 & F06 & Pasien memindai identitas/PRE public key tenaga kesehatan dan memberikan akses \texttt{READ\_PGHD}. & Grant untuk pasangan pasien--tenaga kesehatan tercatat dan material akses sementara dibentuk. \\
    \hline
    T-F07 & F07 & Tenaga kesehatan dengan grant aktif membuka daftar PGHD pasien. & Daftar metadata PGHD valid milik pasien yang sesuai ditampilkan. \\
    \hline
    T-F08 & F08 & Tenaga kesehatan memilih satu entri pada daftar dan membuka detail. & Data berhasil diambil, diverifikasi, didekripsi, dan ditampilkan bersama metadata provenance. \\
    \hline
    T-F09 & F09 & Tenaga kesehatan mencoba membuka fixture PGHD dengan hash, tanda tangan, atau plaintext yang telah diubah. & Sistem menolak data, menampilkan peringatan integritas, dan tidak memperlakukannya sebagai PGHD valid. \\
    \hline
    T-F10 & F10 & Pasien mencabut grant, lalu tenaga kesehatan mencoba menyegarkan daftar dan membuka detail kembali. & Permintaan baca berikutnya ditolak dan tidak menghasilkan plaintext atau material dekripsi. \\
    \hline
\end{styledxltabular}

Skenario T-F01--T-F10 berfungsi sebagai pengujian sistem dari perspektif pengguna. Logika negatif yang deterministik, seperti tipe akses salah, pasien berbeda, akses kedaluwarsa, dan status data invalid, diuji lebih rinci pada pengujian \textit{smart contract}.
